Public fresh/rotten produce datasets are routinely reported with accuracies
close to 100\%. We show that, for a widely used Kaggle dataset, much of this
figure is an artefact of how the data are split and how the data were
collected. Its \NImages{} images derive from only \NGroups{} source
photographs taken in \NSessions{} (file-name-inferred) capture sessions. A per-file split places a
sibling of \NaiveLeakPct\% of test images in training, and an EfficientNet-B0
then scores \NaiveFreshAcc\% on freshness. Grouping copies by source photo
gives \MtlFreshAcc\%, yet \SessionSharePct\% of test images still share a
capture session with training. Holding out whole sessions lowers freshness
accuracy to \SessFreshAcc\%, and expected calibration error grows from
\MtlFreshEce\% to \SessFreshEce\%. The loss is concentrated in produce types
whose fresh and stale photographs came from different cameras or messaging
apps: on an unseen session of stale capsicum the model is correct on only
\SessHeldCapsicumStale\% of images, suggesting it has partly learned the
acquisition source rather than freshness. On the photo-grouped split, joint
training of freshness and produce type shows no difference from separate
single-task models that three seeds can resolve, and MobileNetV3-Large matches EfficientNet-B0 (\MnvFreshAcc\% vs.
\MtlFreshAcc\%) at \MnvCpuMs{}\,ms CPU latency. An energy-score gate on the
produce-type head rejects \MnvGateFarOODRejectRate\% of non-produce images and
\MnvGateNearOODRejectRate\% of unseen produce while accepting
\MnvGateTestIdAcceptRate\% of in-distribution images; produce-type accuracy on
an external dataset is only \MnvCrossTypeAcc\%. We release splits and code, and
argue that shelf life and quality grade, which have no ground truth in these
datasets, should not be presented as learned outputs.
